\documentclass[aspectratio=169,11pt]{beamer}

\usetheme{Madrid}
\usecolortheme{default}
\usefonttheme{professionalfonts}
\setbeamertemplate{navigation symbols}{}
\setbeamertemplate{footline}[frame number]

\usepackage{amsmath, amssymb, booktabs, graphicx}

\newcommand{\E}{\mathbb{E}}
\newcommand{\1}{\mathbf{1}}
\newcommand{\argmax}{\operatorname*{arg\,max}}
\newcommand{\Prb}{\operatorname{Pr}}

\title{Causal ML, Policy Learning, and External Validity}
\subtitle{Empirical Methods --- Week 11}
\author{Teaching deck draft}
\date{}

\begin{document}

\begin{frame}
  \titlepage
\end{frame}

\begin{frame}{Week 11 question}
\begin{itemize}
    \item What should researchers do after heterogeneity is estimated?
    \item When is heterogeneity itself the research object?
    \item When is the object a policy rule, value, or regret comparison?
    \item When does a rule learned in one setting fail to travel to another?
\end{itemize}

\medskip
The method is useful only after the welfare objective, feasible rule class, and deployment environment are explicit.
\end{frame}

\begin{frame}{Lecture 10 versus Lecture 11}
\small
\begin{tabular}{p{0.45\linewidth} p{0.45\linewidth}}
\toprule
Lecture 10 & Lecture 11 \\
\midrule
High-dimensional controls & Policy rules and targeting \\
Orthogonal nuisance estimation & Policy value and regret \\
DML for low-dimensional effects & Empirical welfare maximization \\
CATE and causal forests & CATEs as inputs to decisions \\
Overlap for effect estimation & Overlap where the rule assigns treatment \\
Inference for effects & Validation for value, portability, and fairness \\
\bottomrule
\end{tabular}

\medskip
The continuity is real, but the target object changes.
\end{frame}

\begin{frame}{Three objects that often get confused}
\small
\begin{tabular}{p{0.25\linewidth} p{0.34\linewidth} p{0.31\linewidth}}
\toprule
Object & Mathematical object & Applied question \\
\midrule
Prediction & \(\widehat m(x)\approx \E[Y\mid X=x]\) & What outcome or risk do we expect? \\
Causal effect & \(\tau(x)=\E[Y(1)-Y(0)\mid X=x]\) & For whom does treatment change outcomes? \\
Policy rule & \(\pi(x)\in\{0,1\}\) & Who should be assigned, under constraints? \\
\bottomrule
\end{tabular}

\medskip
High predicted outcomes are not the same as high treatment effects.
\end{frame}

\begin{frame}{CATE versus policy rule}
\small
\[
\tau_i=Y_i(1)-Y_i(0),
\qquad
\tau_G=\E[Y(1)-Y(0)\mid X\in G]
\]

\[
\tau(x)=\E[Y(1)-Y(0)\mid X=x],
\qquad
\pi(x)\in\{0,1\}
\]

\begin{itemize}
    \item CATEs summarize heterogeneity under an identifying design.
    \item Subgroup effects answer incidence or mechanism questions.
    \item A policy rule adds welfare weights, costs, capacity, legality, information, and implementation.
    \item The simple rule \(\pi^*(x)=\1\{\tau(x)\geq 0\}\) needs costless treatment and no constraints.
\end{itemize}
\end{frame}

\begin{frame}{Policy value and regret}
\small
\[
V(\pi)=\E[Y(\pi(X))]
=
\E\left[\pi(X)Y(1)+\left(1-\pi(X)\right)Y(0)\right]
\]

\[
V^{net}(\pi)
=
\E\left[\pi(X)(Y(1)-c(X))+\left(1-\pi(X)\right)Y(0)\right]
\]

\[
R(\pi)=V(\pi^*)-V(\pi)
\]

\begin{itemize}
    \item Report value relative to treat-all, treat-none, simple rules, and current assignment.
    \item Evaluate value on held-out or cross-fit scores, not only in the sample used to choose the rule.
    \item A stable simple rule can have low regret even when a flexible rule has higher in-sample value.
\end{itemize}
\end{frame}

\begin{frame}{Empirical welfare maximization}
\small
\[
\widehat\pi \in \argmax_{\pi\in\Pi}\widehat V(\pi)
\]

\begin{itemize}
    \item \(\Pi\) is the feasible policy class: thresholds, trees, score rules, or constrained assignment rules.
    \item \(\widehat V(\pi)\) estimates welfare or net outcomes under the rule.
    \item Capacity, budget, and fairness constraints belong in \(\Pi\), not in an afterthought paragraph.
    \item Kitagawa--Tetenov and Athey--Wager treat rules as empirical objects whose value can be estimated.
\end{itemize}

\medskip
Policy learning asks which rule performs best, not which subgroup has the largest noisy effect estimate.
\end{frame}

\begin{frame}{Doubly robust policy value}
\small
\[
\mu_d(x)=\E[Y\mid D=d,X=x],
\qquad
e(x)=\Prb(D=1\mid X=x)
\]

\[
\Gamma_i(\pi)=
\pi(X_i)\left[\widehat\mu_1(X_i)+
\frac{D_i}{\widehat e(X_i)}(Y_i-\widehat\mu_1(X_i))\right]
\]
\[
+(1-\pi(X_i))\left[\widehat\mu_0(X_i)+
\frac{1-D_i}{1-\widehat e(X_i)}(Y_i-\widehat\mu_0(X_i))\right]
\]

\[
\widehat V_{DR}(\pi)=\frac{1}{n}\sum_i\Gamma_i(\pi)
\]

\medskip
This inherits Lecture 10's discipline: identify first, learn nuisance functions honestly, and evaluate the target on held-out scores.
\end{frame}

\begin{frame}{External validity and transportability}
\small
\[
V_T(\pi)=\E_T[Y(\pi(X))]
\]

\[
\E_T[h(X)] = \E_S[\omega(X)h(X)],
\qquad
\omega(X)=\frac{f_T(X)}{f_S(X)}
\]

\[
V_T(\pi)\approx
\E_S\left[\omega(X)\{\mu_0(X)+\pi(X)\tau(X)\}\right]
\]

\begin{itemize}
    \item Reweighting helps only when relevant state variables are observed and support overlaps.
    \item Target populations may differ in effects, take-up, costs, implementation, prices, or equilibrium response.
    \item External validity is a design claim, not a property of an algorithm.
\end{itemize}
\end{frame}

\begin{frame}{Fairness and implementation constraints}
\small
\begin{itemize}
    \item Is every feature observable before assignment?
    \item Is the rule legal and administratively feasible?
    \item Does assignment create disparate impact across protected groups?
    \item Are protected groups proxied by geography, language, school, firm, or case history?
    \item Can workers, firms, schools, or staff game the rule?
    \item Does targeting change take-up, stigma, search, or local equilibrium behavior?
\end{itemize}

\medskip
Statistical optimality does not settle legitimacy, feasibility, or welfare.
\end{frame}

\begin{frame}{ML algorithm quick guide}
\scriptsize
\begin{tabular}{p{0.23\linewidth} p{0.34\linewidth} p{0.33\linewidth}}
\toprule
Tool & Most suitable for & Main applied caveat \\
\midrule
Lasso / elastic net & high-dimensional controls, sparse prediction & selected variables are not mechanisms \\
Ridge & dense correlated predictors & weak interpretability \\
Random forests / boosting & nonlinear tabular prediction & feature importance is not causal importance \\
Causal forests / GRF & heterogeneous effects under a design & CATEs are not rules by themselves \\
Policy trees & interpretable assignment rules & unstable when samples are small \\
SVM & margin-based classification & not a natural welfare or causal tool \\
Neural nets & text, images, sequences, very large data & opaque and often excessive for tabular economics \\
Clustering & exploratory typologies & clusters are not causal groups \\
\bottomrule
\end{tabular}
\end{frame}

\begin{frame}{Research Lab design}
\begin{columns}[T,onlytextwidth]
\begin{column}{0.32\linewidth}
\textbf{Reproduce}
\begin{itemize}
    \item synthetic job-training data
    \item CATE by subgroup
    \item feasible policy class
    \item value and regret
\end{itemize}
\end{column}
\begin{column}{0.32\linewidth}
\textbf{Diagnose}
\begin{itemize}
    \item held-out value
    \item overlap in targeted group
    \item capacity sensitivity
    \item fairness audit
\end{itemize}
\end{column}
\begin{column}{0.32\linewidth}
\textbf{Transfer}
\begin{itemize}
    \item shifted target population
    \item source-target weights
    \item target value comparison
    \item portability memo
\end{itemize}
\end{column}
\end{columns}

\medskip
The lab reproduces policy-learning logic, not official paper estimates.
\end{frame}

\begin{frame}{Takeaways}
\begin{itemize}
    \item Heterogeneity is not automatically a policy.
    \item Policy learning needs a rule class, value function, costs, constraints, and honest evaluation.
    \item Targeting should be based on causal gains, not predicted outcome levels.
    \item External validity gets harder when effects and institutions vary.
    \item Fairness and implementation are part of applied economics, not decoration after the estimator.
    \item The best empirical recommendation may be a modest rule with clear value and low regret.
\end{itemize}
\end{frame}

\end{document}
